% ============================================================================
% CV template: Minimal violet (the owner's standard CV design, from Work Finder)
% Layout only. Every name, contact, headline and bullet lives in the content
% file, which atelier-tex generates from a CvDocument for each render.
%
% Compile with XeLaTeX, twice:
%   xelatex "\def\ContentFile{content}\def\FontDir{fonts/}\def\CvFont{plex}% ============================================================================
% CV template: Minimal violet (the owner's standard CV design, from Work Finder)
% Layout only. Every name, contact, headline and bullet lives in the content
% file, which atelier-tex generates from a CvDocument for each render.
%
% Compile with XeLaTeX, twice:
%   xelatex "\def\ContentFile{content}\def\FontDir{fonts/}\def\CvFont{plex}% ============================================================================
% CV template: Minimal violet (the owner's standard CV design, from Work Finder)
% Layout only. Every name, contact, headline and bullet lives in the content
% file, which atelier-tex generates from a CvDocument for each render.
%
% Compile with XeLaTeX, twice:
%   xelatex "\def\ContentFile{content}\def\FontDir{fonts/}\def\CvFont{plex}% ============================================================================
% CV template: Minimal violet (the owner's standard CV design, from Work Finder)
% Layout only. Every name, contact, headline and bullet lives in the content
% file, which atelier-tex generates from a CvDocument for each render.
%
% Compile with XeLaTeX, twice:
%   xelatex "\def\ContentFile{content}\def\FontDir{fonts/}\def\CvFont{plex}\input{cv}"
% \ContentFile defaults to "content", \FontDir to "fonts/" (the folder of the
% chosen family under apps/atelier-tex/fonts) and \CvFont to "plex".
% atelier-tex does all of this in a temporary folder per render.
% Changes from Work Finder: an empty location or education date drops its
% middle dot instead of leaving it dangling, and a summary with an empty
% second paragraph drops the gap before it.
%
% Font: IBM Plex Sans by default (see \CvFont). Type scale (pt): name 29,
% headline 12.5, job title 10.6, body 9.5 on 13,
% metadata 8.8, section labels 9.2 tracked. Accent violet #6D28D9.
%
% ATS rules this layout keeps (verified with poppler and a drawing-order
% parser): single text column, real text only, no tables, icons, header or
% footer; no ligatures and no hyphenation; Manrope, which maps every glyph to
% plain ASCII (hyphens and parentheses included); tracking at or below
% 0.06 em; skills one category per line; inline items joined by middle dots,
% never by wide gaps; nothing right-aligned alone where a parser would read
% it as a separate column.
% ============================================================================
\documentclass[a4paper,10pt]{article}

\providecommand{\ContentFile}{content}
\providecommand{\FontDir}{fonts/}

% ----------------------------------------------------------------------------
% Page and packages
% ----------------------------------------------------------------------------
\usepackage[a4paper,left=15mm,right=15mm,top=13mm,bottom=14mm]{geometry}
\usepackage{fontspec}
\usepackage[english]{babel}
\usepackage[protrusion=true]{microtype}
\usepackage{enumitem}
\usepackage{titlesec}
\usepackage{needspace}
\usepackage{etoolbox}
\usepackage{xcolor}
\usepackage[hidelinks,unicode]{hyperref}

% ----------------------------------------------------------------------------
% Fonts: plain letters only, so "fi", "fl" and "ff" stay searchable.
% \CvFont picks the family: plex (IBM Plex Sans, bold text in SemiBold; the
% default), manrope or atkinson (Atkinson Hyperlegible). All three map every
% glyph back to plain ASCII. Plex keeps I, l and 1 distinct, so OCR-based
% checkers read it far better than Manrope.
% ----------------------------------------------------------------------------
\providecommand{\CvFont}{plex}
\ifdefstring{\CvFont}{plex}{%
  \def\CvRegular{IBMPlexSans-Regular.ttf}\def\CvBold{IBMPlexSans-SemiBold.ttf}%
  \def\CvHeavy{IBMPlexSans-Bold.ttf}\def\CvNameTrack{-0.5}}{%
\ifdefstring{\CvFont}{atkinson}{%
  \def\CvRegular{AtkinsonHyperlegible-Regular.ttf}\def\CvBold{AtkinsonHyperlegible-Bold.ttf}%
  \def\CvHeavy{AtkinsonHyperlegible-Bold.ttf}\def\CvNameTrack{0}}{%
  \def\CvRegular{Manrope-Regular.ttf}\def\CvBold{Manrope-Bold.ttf}%
  \def\CvHeavy{Manrope-Bold.ttf}\def\CvNameTrack{-1.5}}}
\defaultfontfeatures{Ligatures={TeX,NoCommon},Kerning=On}
\setmainfont{\CvRegular}[Path=\FontDir,BoldFont=\CvBold]
\newfontfamily\namefont{\CvHeavy}[Path=\FontDir,LetterSpace=\CvNameTrack]
% Tracking stays below 0.1 em, so parsers still read every label as one word.
\newfontfamily\labelfont{\CvBold}[Path=\FontDir,LetterSpace=6]

% No hyphenation: every keyword stays whole for ATS parsers.
\hyphenpenalty=10000
\exhyphenpenalty=10000
\setlength{\emergencystretch}{3em}

% ----------------------------------------------------------------------------
% Colors: one violet accent, a deep violet-ink for highlights, neutral text
% ----------------------------------------------------------------------------
\providecommand{\AccentHex}{6D28D9}
\definecolor{Accent}{HTML}{\AccentHex}
\colorlet{Highlight}{Accent!32!black}
\colorlet{Rule}{Accent!18!white}
\colorlet{Divider}{Accent!45!white}
\definecolor{Ink}{HTML}{14171C}
\definecolor{Body}{HTML}{2B3038}
\definecolor{Muted}{HTML}{646B76}
\setlength{\parindent}{0pt}
\setlength{\parskip}{0pt}
\raggedbottom
% Never split a bullet so that one of its lines sits alone on a page.
\clubpenalty=10000
\widowpenalty=10000
\displaywidowpenalty=10000

% No header or footer: running text there lands in the middle of a job's
% bullets when ATS parsers extract the page.
\pagestyle{empty}

% ----------------------------------------------------------------------------
% Header in three rows: the name with location and phone at the right edge,
% the headline with the core stack at the right edge, then every link on one
% line. Location and phone sit a few points off the name's baseline, so
% parsers keep the name as a line of its own. Links use the accent so they
% read as clickable.
%   \cvheader{name}{location}{phone as shown}{phone digits}{headline}{stack}
%   \cvlinks{\cvlink{url}{text}\linksep\cvlink{url}{text}...}
% ----------------------------------------------------------------------------
\newcommand{\linksep}{\hspace{7pt}{\color{Divider}·}\hspace{7pt}}
\newcommand{\cvlink}[2]{\href{#1}{\color{Accent}#2}}
\newcommand{\cvheader}[6]{%
  \hypersetup{pdfauthor={#1}}%
  \noindent{\namefont\fontsize{29}{33}\selectfont\color{Ink}#1}\hfill
  % Raised only slightly: the block's top must stay below the top of the name,
  % or poppler reads location and phone before the name.
  \raisebox{0.5pt}{\parbox[b]{0.4\textwidth}{\raggedleft\fontsize{9.5}{12.5}\selectfont
    {\bfseries\color{Ink}#2}\par
    {\color{Ink}\href{tel:#4}{#3}}}}\par
  \vspace{3pt}
  {\fontsize{12.5}{16}\selectfont\bfseries\color{Accent}#5}\hfill
  {\fontsize{9.5}{16}\selectfont\color{Muted}#6}\par
  \vspace{7pt}}
\newcommand{\cvlinks}[1]{{\fontsize{9.2}{12}\selectfont#1\par}\vspace{2pt}}

% ----------------------------------------------------------------------------
% Sections and blocks
% ----------------------------------------------------------------------------
% Section: tracked uppercase label followed by a hairline to the margin.
\newcommand{\sectionlabel}[1]{\MakeUppercase{#1}\hspace{10pt}{\color{Rule}\leaders\vrule height 3.5pt depth -2.85pt\hfill}\kern0pt}
\titleformat{\section}{\labelfont\fontsize{9.2}{12}\selectfont\color{Accent}}{}{0pt}{\sectionlabel}
\titlespacing*{\section}{0pt}{10pt}{5.5pt}
\newcommand{\cvsection}[1]{\Needspace{4\baselineskip}\section*{#1}}

\newcommand{\hl}[1]{{\bfseries\color{Highlight}#1}}

\newcommand{\cvsummary}[2]{#1\par\ifblank{#2}{}{\vspace{4.5pt}#2\par}}

% Skills: bold label and the list on the same line, one line per category.
\newenvironment{cvskills}{\par}{\par}
\newcommand{\cvskill}[2]{%
  \par\hangindent=10pt\hangafter=1\noindent
  {\bfseries\color{Ink}#1:}\hspace{4pt}#2\par\vspace{2.6pt}}

% Job: title | company on the left, dates and location on the right,
% a one-line company description underneath; never stranded at a page end.
%   \cvjob{title}{company}{dates}{location}{company description}
\newcommand{\cvjob}[5]{%
  \par\addvspace{8pt}\Needspace{4\baselineskip}%
  {\fontsize{10.6}{14}\selectfont{\bfseries\color{Ink}#1}\hspace{7pt}{\color{Divider}|}\hspace{7pt}{\bfseries\color{Accent}#2}\hfill{\fontsize{8.8}{14}\selectfont\color{Muted}#3\ifblank{#4}{}{\hspace{5pt}·\hspace{5pt}#4}}}\par
  {\fontsize{8.8}{12.4}\selectfont\color{Muted}#5}\par\nopagebreak
  \vspace{3.5pt}}

\newenvironment{cvbullets}{%
  \begin{itemize}[leftmargin=12.5pt,labelsep=6.5pt,itemsep=2pt,topsep=0pt,parsep=0pt,partopsep=0pt,
    label={\color{Accent}\fontsize{8.5}{8.5}\selectfont\textbullet}]}
  {\end{itemize}}

% Education: degree, then school and dates on one left-aligned line. Dates
% pushed to the right margin here would be read as a separate column.
\newcommand{\cvedu}[3]{%
  {\fontsize{10.6}{14}\selectfont\bfseries\color{Ink}#1}\par
  {\bfseries\color{Accent}#2}\ifblank{#3}{}{\hspace{6pt}{\color{Muted}·\hspace{6pt}#3}}\par}

% Languages share one line, joined by a middle dot. A wide gap instead would
% make parsers read the second language as a separate column.
\newcommand{\cvlangsep}{}
\newcommand{\cvlang}[2]{\cvlangsep{\bfseries\color{Ink}#1:}\hspace{4pt}#2%
  \gdef\cvlangsep{\hspace{6pt}{\color{Muted}·}\hspace{6pt}}}

% ============================================================================
\begin{document}
\fontsize{9.5}{13}\selectfont
\color{Body}
\raggedright
\frenchspacing

\input{\ContentFile}

\end{document}
"
% \ContentFile defaults to "content", \FontDir to "fonts/" (the folder of the
% chosen family under apps/atelier-tex/fonts) and \CvFont to "plex".
% atelier-tex does all of this in a temporary folder per render.
% Changes from Work Finder: an empty location or education date drops its
% middle dot instead of leaving it dangling, and a summary with an empty
% second paragraph drops the gap before it.
%
% Font: IBM Plex Sans by default (see \CvFont). Type scale (pt): name 29,
% headline 12.5, job title 10.6, body 9.5 on 13,
% metadata 8.8, section labels 9.2 tracked. Accent violet #6D28D9.
%
% ATS rules this layout keeps (verified with poppler and a drawing-order
% parser): single text column, real text only, no tables, icons, header or
% footer; no ligatures and no hyphenation; Manrope, which maps every glyph to
% plain ASCII (hyphens and parentheses included); tracking at or below
% 0.06 em; skills one category per line; inline items joined by middle dots,
% never by wide gaps; nothing right-aligned alone where a parser would read
% it as a separate column.
% ============================================================================
\documentclass[a4paper,10pt]{article}

\providecommand{\ContentFile}{content}
\providecommand{\FontDir}{fonts/}

% ----------------------------------------------------------------------------
% Page and packages
% ----------------------------------------------------------------------------
\usepackage[a4paper,left=15mm,right=15mm,top=13mm,bottom=14mm]{geometry}
\usepackage{fontspec}
\usepackage[english]{babel}
\usepackage[protrusion=true]{microtype}
\usepackage{enumitem}
\usepackage{titlesec}
\usepackage{needspace}
\usepackage{etoolbox}
\usepackage{xcolor}
\usepackage[hidelinks,unicode]{hyperref}

% ----------------------------------------------------------------------------
% Fonts: plain letters only, so "fi", "fl" and "ff" stay searchable.
% \CvFont picks the family: plex (IBM Plex Sans, bold text in SemiBold; the
% default), manrope or atkinson (Atkinson Hyperlegible). All three map every
% glyph back to plain ASCII. Plex keeps I, l and 1 distinct, so OCR-based
% checkers read it far better than Manrope.
% ----------------------------------------------------------------------------
\providecommand{\CvFont}{plex}
\ifdefstring{\CvFont}{plex}{%
  \def\CvRegular{IBMPlexSans-Regular.ttf}\def\CvBold{IBMPlexSans-SemiBold.ttf}%
  \def\CvHeavy{IBMPlexSans-Bold.ttf}\def\CvNameTrack{-0.5}}{%
\ifdefstring{\CvFont}{atkinson}{%
  \def\CvRegular{AtkinsonHyperlegible-Regular.ttf}\def\CvBold{AtkinsonHyperlegible-Bold.ttf}%
  \def\CvHeavy{AtkinsonHyperlegible-Bold.ttf}\def\CvNameTrack{0}}{%
  \def\CvRegular{Manrope-Regular.ttf}\def\CvBold{Manrope-Bold.ttf}%
  \def\CvHeavy{Manrope-Bold.ttf}\def\CvNameTrack{-1.5}}}
\defaultfontfeatures{Ligatures={TeX,NoCommon},Kerning=On}
\setmainfont{\CvRegular}[Path=\FontDir,BoldFont=\CvBold]
\newfontfamily\namefont{\CvHeavy}[Path=\FontDir,LetterSpace=\CvNameTrack]
% Tracking stays below 0.1 em, so parsers still read every label as one word.
\newfontfamily\labelfont{\CvBold}[Path=\FontDir,LetterSpace=6]

% No hyphenation: every keyword stays whole for ATS parsers.
\hyphenpenalty=10000
\exhyphenpenalty=10000
\setlength{\emergencystretch}{3em}

% ----------------------------------------------------------------------------
% Colors: one violet accent, a deep violet-ink for highlights, neutral text
% ----------------------------------------------------------------------------
\providecommand{\AccentHex}{6D28D9}
\definecolor{Accent}{HTML}{\AccentHex}
\colorlet{Highlight}{Accent!32!black}
\colorlet{Rule}{Accent!18!white}
\colorlet{Divider}{Accent!45!white}
\definecolor{Ink}{HTML}{14171C}
\definecolor{Body}{HTML}{2B3038}
\definecolor{Muted}{HTML}{646B76}
\setlength{\parindent}{0pt}
\setlength{\parskip}{0pt}
\raggedbottom
% Never split a bullet so that one of its lines sits alone on a page.
\clubpenalty=10000
\widowpenalty=10000
\displaywidowpenalty=10000

% No header or footer: running text there lands in the middle of a job's
% bullets when ATS parsers extract the page.
\pagestyle{empty}

% ----------------------------------------------------------------------------
% Header in three rows: the name with location and phone at the right edge,
% the headline with the core stack at the right edge, then every link on one
% line. Location and phone sit a few points off the name's baseline, so
% parsers keep the name as a line of its own. Links use the accent so they
% read as clickable.
%   \cvheader{name}{location}{phone as shown}{phone digits}{headline}{stack}
%   \cvlinks{\cvlink{url}{text}\linksep\cvlink{url}{text}...}
% ----------------------------------------------------------------------------
\newcommand{\linksep}{\hspace{7pt}{\color{Divider}·}\hspace{7pt}}
\newcommand{\cvlink}[2]{\href{#1}{\color{Accent}#2}}
\newcommand{\cvheader}[6]{%
  \hypersetup{pdfauthor={#1}}%
  \noindent{\namefont\fontsize{29}{33}\selectfont\color{Ink}#1}\hfill
  % Raised only slightly: the block's top must stay below the top of the name,
  % or poppler reads location and phone before the name.
  \raisebox{0.5pt}{\parbox[b]{0.4\textwidth}{\raggedleft\fontsize{9.5}{12.5}\selectfont
    {\bfseries\color{Ink}#2}\par
    {\color{Ink}\href{tel:#4}{#3}}}}\par
  \vspace{3pt}
  {\fontsize{12.5}{16}\selectfont\bfseries\color{Accent}#5}\hfill
  {\fontsize{9.5}{16}\selectfont\color{Muted}#6}\par
  \vspace{7pt}}
\newcommand{\cvlinks}[1]{{\fontsize{9.2}{12}\selectfont#1\par}\vspace{2pt}}

% ----------------------------------------------------------------------------
% Sections and blocks
% ----------------------------------------------------------------------------
% Section: tracked uppercase label followed by a hairline to the margin.
\newcommand{\sectionlabel}[1]{\MakeUppercase{#1}\hspace{10pt}{\color{Rule}\leaders\vrule height 3.5pt depth -2.85pt\hfill}\kern0pt}
\titleformat{\section}{\labelfont\fontsize{9.2}{12}\selectfont\color{Accent}}{}{0pt}{\sectionlabel}
\titlespacing*{\section}{0pt}{10pt}{5.5pt}
\newcommand{\cvsection}[1]{\Needspace{4\baselineskip}\section*{#1}}

\newcommand{\hl}[1]{{\bfseries\color{Highlight}#1}}

\newcommand{\cvsummary}[2]{#1\par\ifblank{#2}{}{\vspace{4.5pt}#2\par}}

% Skills: bold label and the list on the same line, one line per category.
\newenvironment{cvskills}{\par}{\par}
\newcommand{\cvskill}[2]{%
  \par\hangindent=10pt\hangafter=1\noindent
  {\bfseries\color{Ink}#1:}\hspace{4pt}#2\par\vspace{2.6pt}}

% Job: title | company on the left, dates and location on the right,
% a one-line company description underneath; never stranded at a page end.
%   \cvjob{title}{company}{dates}{location}{company description}
\newcommand{\cvjob}[5]{%
  \par\addvspace{8pt}\Needspace{4\baselineskip}%
  {\fontsize{10.6}{14}\selectfont{\bfseries\color{Ink}#1}\hspace{7pt}{\color{Divider}|}\hspace{7pt}{\bfseries\color{Accent}#2}\hfill{\fontsize{8.8}{14}\selectfont\color{Muted}#3\ifblank{#4}{}{\hspace{5pt}·\hspace{5pt}#4}}}\par
  {\fontsize{8.8}{12.4}\selectfont\color{Muted}#5}\par\nopagebreak
  \vspace{3.5pt}}

\newenvironment{cvbullets}{%
  \begin{itemize}[leftmargin=12.5pt,labelsep=6.5pt,itemsep=2pt,topsep=0pt,parsep=0pt,partopsep=0pt,
    label={\color{Accent}\fontsize{8.5}{8.5}\selectfont\textbullet}]}
  {\end{itemize}}

% Education: degree, then school and dates on one left-aligned line. Dates
% pushed to the right margin here would be read as a separate column.
\newcommand{\cvedu}[3]{%
  {\fontsize{10.6}{14}\selectfont\bfseries\color{Ink}#1}\par
  {\bfseries\color{Accent}#2}\ifblank{#3}{}{\hspace{6pt}{\color{Muted}·\hspace{6pt}#3}}\par}

% Languages share one line, joined by a middle dot. A wide gap instead would
% make parsers read the second language as a separate column.
\newcommand{\cvlangsep}{}
\newcommand{\cvlang}[2]{\cvlangsep{\bfseries\color{Ink}#1:}\hspace{4pt}#2%
  \gdef\cvlangsep{\hspace{6pt}{\color{Muted}·}\hspace{6pt}}}

% ============================================================================
\begin{document}
\fontsize{9.5}{13}\selectfont
\color{Body}
\raggedright
\frenchspacing

\input{\ContentFile}

\end{document}
"
% \ContentFile defaults to "content", \FontDir to "fonts/" (the folder of the
% chosen family under apps/atelier-tex/fonts) and \CvFont to "plex".
% atelier-tex does all of this in a temporary folder per render.
% Changes from Work Finder: an empty location or education date drops its
% middle dot instead of leaving it dangling, and a summary with an empty
% second paragraph drops the gap before it.
%
% Font: IBM Plex Sans by default (see \CvFont). Type scale (pt): name 29,
% headline 12.5, job title 10.6, body 9.5 on 13,
% metadata 8.8, section labels 9.2 tracked. Accent violet #6D28D9.
%
% ATS rules this layout keeps (verified with poppler and a drawing-order
% parser): single text column, real text only, no tables, icons, header or
% footer; no ligatures and no hyphenation; Manrope, which maps every glyph to
% plain ASCII (hyphens and parentheses included); tracking at or below
% 0.06 em; skills one category per line; inline items joined by middle dots,
% never by wide gaps; nothing right-aligned alone where a parser would read
% it as a separate column.
% ============================================================================
\documentclass[a4paper,10pt]{article}

\providecommand{\ContentFile}{content}
\providecommand{\FontDir}{fonts/}

% ----------------------------------------------------------------------------
% Page and packages
% ----------------------------------------------------------------------------
\usepackage[a4paper,left=15mm,right=15mm,top=13mm,bottom=14mm]{geometry}
\usepackage{fontspec}
\usepackage[english]{babel}
\usepackage[protrusion=true]{microtype}
\usepackage{enumitem}
\usepackage{titlesec}
\usepackage{needspace}
\usepackage{etoolbox}
\usepackage{xcolor}
\usepackage[hidelinks,unicode]{hyperref}

% ----------------------------------------------------------------------------
% Fonts: plain letters only, so "fi", "fl" and "ff" stay searchable.
% \CvFont picks the family: plex (IBM Plex Sans, bold text in SemiBold; the
% default), manrope or atkinson (Atkinson Hyperlegible). All three map every
% glyph back to plain ASCII. Plex keeps I, l and 1 distinct, so OCR-based
% checkers read it far better than Manrope.
% ----------------------------------------------------------------------------
\providecommand{\CvFont}{plex}
\ifdefstring{\CvFont}{plex}{%
  \def\CvRegular{IBMPlexSans-Regular.ttf}\def\CvBold{IBMPlexSans-SemiBold.ttf}%
  \def\CvHeavy{IBMPlexSans-Bold.ttf}\def\CvNameTrack{-0.5}}{%
\ifdefstring{\CvFont}{atkinson}{%
  \def\CvRegular{AtkinsonHyperlegible-Regular.ttf}\def\CvBold{AtkinsonHyperlegible-Bold.ttf}%
  \def\CvHeavy{AtkinsonHyperlegible-Bold.ttf}\def\CvNameTrack{0}}{%
  \def\CvRegular{Manrope-Regular.ttf}\def\CvBold{Manrope-Bold.ttf}%
  \def\CvHeavy{Manrope-Bold.ttf}\def\CvNameTrack{-1.5}}}
\defaultfontfeatures{Ligatures={TeX,NoCommon},Kerning=On}
\setmainfont{\CvRegular}[Path=\FontDir,BoldFont=\CvBold]
\newfontfamily\namefont{\CvHeavy}[Path=\FontDir,LetterSpace=\CvNameTrack]
% Tracking stays below 0.1 em, so parsers still read every label as one word.
\newfontfamily\labelfont{\CvBold}[Path=\FontDir,LetterSpace=6]

% No hyphenation: every keyword stays whole for ATS parsers.
\hyphenpenalty=10000
\exhyphenpenalty=10000
\setlength{\emergencystretch}{3em}

% ----------------------------------------------------------------------------
% Colors: one violet accent, a deep violet-ink for highlights, neutral text
% ----------------------------------------------------------------------------
\providecommand{\AccentHex}{6D28D9}
\definecolor{Accent}{HTML}{\AccentHex}
\colorlet{Highlight}{Accent!32!black}
\colorlet{Rule}{Accent!18!white}
\colorlet{Divider}{Accent!45!white}
\definecolor{Ink}{HTML}{14171C}
\definecolor{Body}{HTML}{2B3038}
\definecolor{Muted}{HTML}{646B76}
\setlength{\parindent}{0pt}
\setlength{\parskip}{0pt}
\raggedbottom
% Never split a bullet so that one of its lines sits alone on a page.
\clubpenalty=10000
\widowpenalty=10000
\displaywidowpenalty=10000

% No header or footer: running text there lands in the middle of a job's
% bullets when ATS parsers extract the page.
\pagestyle{empty}

% ----------------------------------------------------------------------------
% Header in three rows: the name with location and phone at the right edge,
% the headline with the core stack at the right edge, then every link on one
% line. Location and phone sit a few points off the name's baseline, so
% parsers keep the name as a line of its own. Links use the accent so they
% read as clickable.
%   \cvheader{name}{location}{phone as shown}{phone digits}{headline}{stack}
%   \cvlinks{\cvlink{url}{text}\linksep\cvlink{url}{text}...}
% ----------------------------------------------------------------------------
\newcommand{\linksep}{\hspace{7pt}{\color{Divider}·}\hspace{7pt}}
\newcommand{\cvlink}[2]{\href{#1}{\color{Accent}#2}}
\newcommand{\cvheader}[6]{%
  \hypersetup{pdfauthor={#1}}%
  \noindent{\namefont\fontsize{29}{33}\selectfont\color{Ink}#1}\hfill
  % Raised only slightly: the block's top must stay below the top of the name,
  % or poppler reads location and phone before the name.
  \raisebox{0.5pt}{\parbox[b]{0.4\textwidth}{\raggedleft\fontsize{9.5}{12.5}\selectfont
    {\bfseries\color{Ink}#2}\par
    {\color{Ink}\href{tel:#4}{#3}}}}\par
  \vspace{3pt}
  {\fontsize{12.5}{16}\selectfont\bfseries\color{Accent}#5}\hfill
  {\fontsize{9.5}{16}\selectfont\color{Muted}#6}\par
  \vspace{7pt}}
\newcommand{\cvlinks}[1]{{\fontsize{9.2}{12}\selectfont#1\par}\vspace{2pt}}

% ----------------------------------------------------------------------------
% Sections and blocks
% ----------------------------------------------------------------------------
% Section: tracked uppercase label followed by a hairline to the margin.
\newcommand{\sectionlabel}[1]{\MakeUppercase{#1}\hspace{10pt}{\color{Rule}\leaders\vrule height 3.5pt depth -2.85pt\hfill}\kern0pt}
\titleformat{\section}{\labelfont\fontsize{9.2}{12}\selectfont\color{Accent}}{}{0pt}{\sectionlabel}
\titlespacing*{\section}{0pt}{10pt}{5.5pt}
\newcommand{\cvsection}[1]{\Needspace{4\baselineskip}\section*{#1}}

\newcommand{\hl}[1]{{\bfseries\color{Highlight}#1}}

\newcommand{\cvsummary}[2]{#1\par\ifblank{#2}{}{\vspace{4.5pt}#2\par}}

% Skills: bold label and the list on the same line, one line per category.
\newenvironment{cvskills}{\par}{\par}
\newcommand{\cvskill}[2]{%
  \par\hangindent=10pt\hangafter=1\noindent
  {\bfseries\color{Ink}#1:}\hspace{4pt}#2\par\vspace{2.6pt}}

% Job: title | company on the left, dates and location on the right,
% a one-line company description underneath; never stranded at a page end.
%   \cvjob{title}{company}{dates}{location}{company description}
\newcommand{\cvjob}[5]{%
  \par\addvspace{8pt}\Needspace{4\baselineskip}%
  {\fontsize{10.6}{14}\selectfont{\bfseries\color{Ink}#1}\hspace{7pt}{\color{Divider}|}\hspace{7pt}{\bfseries\color{Accent}#2}\hfill{\fontsize{8.8}{14}\selectfont\color{Muted}#3\ifblank{#4}{}{\hspace{5pt}·\hspace{5pt}#4}}}\par
  {\fontsize{8.8}{12.4}\selectfont\color{Muted}#5}\par\nopagebreak
  \vspace{3.5pt}}

\newenvironment{cvbullets}{%
  \begin{itemize}[leftmargin=12.5pt,labelsep=6.5pt,itemsep=2pt,topsep=0pt,parsep=0pt,partopsep=0pt,
    label={\color{Accent}\fontsize{8.5}{8.5}\selectfont\textbullet}]}
  {\end{itemize}}

% Education: degree, then school and dates on one left-aligned line. Dates
% pushed to the right margin here would be read as a separate column.
\newcommand{\cvedu}[3]{%
  {\fontsize{10.6}{14}\selectfont\bfseries\color{Ink}#1}\par
  {\bfseries\color{Accent}#2}\ifblank{#3}{}{\hspace{6pt}{\color{Muted}·\hspace{6pt}#3}}\par}

% Languages share one line, joined by a middle dot. A wide gap instead would
% make parsers read the second language as a separate column.
\newcommand{\cvlangsep}{}
\newcommand{\cvlang}[2]{\cvlangsep{\bfseries\color{Ink}#1:}\hspace{4pt}#2%
  \gdef\cvlangsep{\hspace{6pt}{\color{Muted}·}\hspace{6pt}}}

% ============================================================================
\begin{document}
\fontsize{9.5}{13}\selectfont
\color{Body}
\raggedright
\frenchspacing

\input{\ContentFile}

\end{document}
"
% \ContentFile defaults to "content", \FontDir to "fonts/" (the folder of the
% chosen family under apps/atelier-tex/fonts) and \CvFont to "plex".
% atelier-tex does all of this in a temporary folder per render.
% Changes from Work Finder: an empty location or education date drops its
% middle dot instead of leaving it dangling, and a summary with an empty
% second paragraph drops the gap before it.
%
% Font: IBM Plex Sans by default (see \CvFont). Type scale (pt): name 29,
% headline 12.5, job title 10.6, body 9.5 on 13,
% metadata 8.8, section labels 9.2 tracked. Accent violet #6D28D9.
%
% ATS rules this layout keeps (verified with poppler and a drawing-order
% parser): single text column, real text only, no tables, icons, header or
% footer; no ligatures and no hyphenation; Manrope, which maps every glyph to
% plain ASCII (hyphens and parentheses included); tracking at or below
% 0.06 em; skills one category per line; inline items joined by middle dots,
% never by wide gaps; nothing right-aligned alone where a parser would read
% it as a separate column.
% ============================================================================
\documentclass[a4paper,10pt]{article}

\providecommand{\ContentFile}{content}
\providecommand{\FontDir}{fonts/}

% ----------------------------------------------------------------------------
% Page and packages
% ----------------------------------------------------------------------------
\usepackage[a4paper,left=15mm,right=15mm,top=13mm,bottom=14mm]{geometry}
\usepackage{fontspec}
\usepackage[english]{babel}
\usepackage[protrusion=true]{microtype}
\usepackage{enumitem}
\usepackage{titlesec}
\usepackage{needspace}
\usepackage{etoolbox}
\usepackage{xcolor}
\usepackage[hidelinks,unicode]{hyperref}

% ----------------------------------------------------------------------------
% Fonts: plain letters only, so "fi", "fl" and "ff" stay searchable.
% \CvFont picks the family: plex (IBM Plex Sans, bold text in SemiBold; the
% default), manrope or atkinson (Atkinson Hyperlegible). All three map every
% glyph back to plain ASCII. Plex keeps I, l and 1 distinct, so OCR-based
% checkers read it far better than Manrope.
% ----------------------------------------------------------------------------
\providecommand{\CvFont}{plex}
\ifdefstring{\CvFont}{plex}{%
  \def\CvRegular{IBMPlexSans-Regular.ttf}\def\CvBold{IBMPlexSans-SemiBold.ttf}%
  \def\CvHeavy{IBMPlexSans-Bold.ttf}\def\CvNameTrack{-0.5}}{%
\ifdefstring{\CvFont}{atkinson}{%
  \def\CvRegular{AtkinsonHyperlegible-Regular.ttf}\def\CvBold{AtkinsonHyperlegible-Bold.ttf}%
  \def\CvHeavy{AtkinsonHyperlegible-Bold.ttf}\def\CvNameTrack{0}}{%
  \def\CvRegular{Manrope-Regular.ttf}\def\CvBold{Manrope-Bold.ttf}%
  \def\CvHeavy{Manrope-Bold.ttf}\def\CvNameTrack{-1.5}}}
\defaultfontfeatures{Ligatures={TeX,NoCommon},Kerning=On}
\setmainfont{\CvRegular}[Path=\FontDir,BoldFont=\CvBold]
\newfontfamily\namefont{\CvHeavy}[Path=\FontDir,LetterSpace=\CvNameTrack]
% Tracking stays below 0.1 em, so parsers still read every label as one word.
\newfontfamily\labelfont{\CvBold}[Path=\FontDir,LetterSpace=6]

% No hyphenation: every keyword stays whole for ATS parsers.
\hyphenpenalty=10000
\exhyphenpenalty=10000
\setlength{\emergencystretch}{3em}

% ----------------------------------------------------------------------------
% Colors: one violet accent, a deep violet-ink for highlights, neutral text
% ----------------------------------------------------------------------------
\providecommand{\AccentHex}{6D28D9}
\definecolor{Accent}{HTML}{\AccentHex}
\colorlet{Highlight}{Accent!32!black}
\colorlet{Rule}{Accent!18!white}
\colorlet{Divider}{Accent!45!white}
\definecolor{Ink}{HTML}{14171C}
\definecolor{Body}{HTML}{2B3038}
\definecolor{Muted}{HTML}{646B76}
\setlength{\parindent}{0pt}
\setlength{\parskip}{0pt}
\raggedbottom
% Never split a bullet so that one of its lines sits alone on a page.
\clubpenalty=10000
\widowpenalty=10000
\displaywidowpenalty=10000

% No header or footer: running text there lands in the middle of a job's
% bullets when ATS parsers extract the page.
\pagestyle{empty}

% ----------------------------------------------------------------------------
% Header in three rows: the name with location and phone at the right edge,
% the headline with the core stack at the right edge, then every link on one
% line. Location and phone sit a few points off the name's baseline, so
% parsers keep the name as a line of its own. Links use the accent so they
% read as clickable.
%   \cvheader{name}{location}{phone as shown}{phone digits}{headline}{stack}
%   \cvlinks{\cvlink{url}{text}\linksep\cvlink{url}{text}...}
% ----------------------------------------------------------------------------
\newcommand{\linksep}{\hspace{7pt}{\color{Divider}·}\hspace{7pt}}
\newcommand{\cvlink}[2]{\href{#1}{\color{Accent}#2}}
\newcommand{\cvheader}[6]{%
  \hypersetup{pdfauthor={#1}}%
  \noindent{\namefont\fontsize{29}{33}\selectfont\color{Ink}#1}\hfill
  % Raised only slightly: the block's top must stay below the top of the name,
  % or poppler reads location and phone before the name.
  \raisebox{0.5pt}{\parbox[b]{0.4\textwidth}{\raggedleft\fontsize{9.5}{12.5}\selectfont
    {\bfseries\color{Ink}#2}\par
    {\color{Ink}\href{tel:#4}{#3}}}}\par
  \vspace{3pt}
  {\fontsize{12.5}{16}\selectfont\bfseries\color{Accent}#5}\hfill
  {\fontsize{9.5}{16}\selectfont\color{Muted}#6}\par
  \vspace{7pt}}
\newcommand{\cvlinks}[1]{{\fontsize{9.2}{12}\selectfont#1\par}\vspace{2pt}}

% ----------------------------------------------------------------------------
% Sections and blocks
% ----------------------------------------------------------------------------
% Section: tracked uppercase label followed by a hairline to the margin.
\newcommand{\sectionlabel}[1]{\MakeUppercase{#1}\hspace{10pt}{\color{Rule}\leaders\vrule height 3.5pt depth -2.85pt\hfill}\kern0pt}
\titleformat{\section}{\labelfont\fontsize{9.2}{12}\selectfont\color{Accent}}{}{0pt}{\sectionlabel}
\titlespacing*{\section}{0pt}{10pt}{5.5pt}
\newcommand{\cvsection}[1]{\Needspace{4\baselineskip}\section*{#1}}

\newcommand{\hl}[1]{{\bfseries\color{Highlight}#1}}

\newcommand{\cvsummary}[2]{#1\par\ifblank{#2}{}{\vspace{4.5pt}#2\par}}

% Skills: bold label and the list on the same line, one line per category.
\newenvironment{cvskills}{\par}{\par}
\newcommand{\cvskill}[2]{%
  \par\hangindent=10pt\hangafter=1\noindent
  {\bfseries\color{Ink}#1:}\hspace{4pt}#2\par\vspace{2.6pt}}

% Job: title | company on the left, dates and location on the right,
% a one-line company description underneath; never stranded at a page end.
%   \cvjob{title}{company}{dates}{location}{company description}
\newcommand{\cvjob}[5]{%
  \par\addvspace{8pt}\Needspace{4\baselineskip}%
  {\fontsize{10.6}{14}\selectfont{\bfseries\color{Ink}#1}\hspace{7pt}{\color{Divider}|}\hspace{7pt}{\bfseries\color{Accent}#2}\hfill{\fontsize{8.8}{14}\selectfont\color{Muted}#3\ifblank{#4}{}{\hspace{5pt}·\hspace{5pt}#4}}}\par
  {\fontsize{8.8}{12.4}\selectfont\color{Muted}#5}\par\nopagebreak
  \vspace{3.5pt}}

\newenvironment{cvbullets}{%
  \begin{itemize}[leftmargin=12.5pt,labelsep=6.5pt,itemsep=2pt,topsep=0pt,parsep=0pt,partopsep=0pt,
    label={\color{Accent}\fontsize{8.5}{8.5}\selectfont\textbullet}]}
  {\end{itemize}}

% Education: degree, then school and dates on one left-aligned line. Dates
% pushed to the right margin here would be read as a separate column.
\newcommand{\cvedu}[3]{%
  {\fontsize{10.6}{14}\selectfont\bfseries\color{Ink}#1}\par
  {\bfseries\color{Accent}#2}\ifblank{#3}{}{\hspace{6pt}{\color{Muted}·\hspace{6pt}#3}}\par}

% Languages share one line, joined by a middle dot. A wide gap instead would
% make parsers read the second language as a separate column.
\newcommand{\cvlangsep}{}
\newcommand{\cvlang}[2]{\cvlangsep{\bfseries\color{Ink}#1:}\hspace{4pt}#2%
  \gdef\cvlangsep{\hspace{6pt}{\color{Muted}·}\hspace{6pt}}}

% ============================================================================
\begin{document}
\fontsize{9.5}{13}\selectfont
\color{Body}
\raggedright
\frenchspacing

\input{\ContentFile}

\end{document}
